\chapter{Đóng góp, triết lý khoa học và các mối đe dọa}
\label{chap:dong-gop}

\section{Các đóng góp khoa học}
\label{sec:dong-gop}

Mỗi đóng góp dưới đây được trình bày dưới dạng: khẳng định, bằng chứng (tệp dữ liệu trong kho mã), điều kiện sẽ bác bỏ
nó, và giới hạn.

\subsection{Đo phần ``grounded'' và phần ``guessed''}

\textbf{Khẳng định.} Ở hạng 1, phần lớn độ chính xác của SemLoc đến từ dòng LLM tự đoán, không từ phổ vi phạm ràng buộc.
Trên lượt chạy của nhóm, Acc@1 \DecSemlocRepairedAllAccOneFull\% = \DecSemlocRepairedAllAccOnePhiGuess{} (dòng đoán) +
\DecSemlocRepairedAllAccOnePhiSpectrum{} (phổ); trên nhật ký của tác giả,
\PubSemlocClaudeTZeroEightSemlocOriginalAllAccOneFull\% = \PubSemlocClaudeTZeroEightSemlocOriginalAllAccOnePhiGuess{} +
\PubSemlocClaudeTZeroEightSemlocOriginalAllAccOnePhiSpectrum; trên lỗi thật, Acc@1 của SemLoc bằng đúng Acc@1 của riêng
dòng đoán (\BipAllLlmGuessAccOneMean\%). \textbf{Bằng chứng:} \texttt{data/rq\_decomposition.csv},
\texttt{data/rq3\_bugsinpy.csv}. \textbf{Bác bỏ nếu:} bỏ điểm thưởng mà Acc@1 không giảm; ở đây Acc@1 giảm
\RqTwoRepairedAllSemlocNoPriorDAccOne{} điểm [\RqTwoRepairedAllSemlocNoPriorDAccOneLo;
\RqTwoRepairedAllSemlocNoPriorDAccOneHi]. \textbf{Giới hạn:} giá trị Shapley phụ thuộc cách định nghĩa ``tắt'' một tín
hiệu; dòng đoán thường lệch một dòng so với nhãn (Chương~\ref{chap:that-bai}).

\subsection{Với lỗi được kiểm thử chạy tới, một prompt có bằng chứng thực thi vượt quy trình nhiều bước}

\textbf{Khẳng định.} Trên bộ dòng lỗi đã sửa (các lỗi mà một kiểm thử thất bại chạy tới), phương pháp đối chứng LLM + độ
phủ có Acc@3 \RqTwoRepairedAllLlmCoverageAccThreeMean\% so với
\RqTwoRepairedAllSemlocAccThreeMean\% của SemLoc ($\Delta$ \RqTwoRepairedAllLlmCoverageDAccThree{}
[\RqTwoRepairedAllLlmCoverageDAccThreeLo; \RqTwoRepairedAllLlmCoverageDAccThreeHi]) và Acc@1 cao hơn
($\Delta$ \RqTwoRepairedAllLlmCoverageDAccOne{} [\RqTwoRepairedAllLlmCoverageDAccOneLo;
\RqTwoRepairedAllLlmCoverageDAccOneHi]); vẫn hơn trên lỗi ngữ nghĩa thuần; tốn \CostLlmCoverageTokensK{} nghìn token mỗi
chương trình so với \CostSemlocTokensK{} nghìn; gần như không đổi khi đổi tên định danh ($\Delta$Acc@3
\RenRepairedLlmCoverageDAccThree). \textbf{Bằng chứng:} \texttt{data/rq2\_ours.csv}, \texttt{data/rq3\_renamed.csv},
\texttt{data/llm\_cost.csv}. \textbf{Giới hạn:} prompt khác chế độ trực tiếp của SemLoc ở nhiều điểm cùng lúc (độ phủ,
kết quả kiểm thử, giả định một dòng lỗi), nên chênh lệch không quy riêng cho độ phủ; trên toàn bộ nhãn gốc, kể cả lỗi
không kiểm thử nào chạy tới, phần hơn nhỏ hơn nhiều và SemLoc tốt hơn ở Acc@5; trên 31 lỗi thật, chênh lệch cùng chiều nhưng khoảng tin cậy còn chứa 0;
so sánh theo từng mô hình chưa có khoảng tin cậy; prompt do nhóm thiết kế.

\subsection{Kết quả công bố phụ thuộc vào cách chấm và thành phần không được mô tả}

\textbf{Khẳng định.} Tái hiện khớp toàn bộ cấu hình, nhưng: điểm thưởng ẩn cho dòng đoán chiếm phần lớn Acc@1; gỡ lỗi chia
đôi đạt \RqOneDdAuthorAllAccOne\% theo mã chấm của tác giả và \RqOneDdFairAllAccOne\% khi chấm công bằng; Acc@5 của SBFL
bị đánh giá thấp (\RqOneSbflOchiaiAuthorAllAccFive\% so với \RqOneSbflOchiaiFairAllAccFive\%); bước kiểm chứng phản
thực tế không đổi thứ hạng nào dù chiếm \CostCfSharePct\% lượt gọi LLM; bộ chạy BugsInPy đo hạng của dòng đoán thay
vì dòng đã sửa và để lại bản biên dịch cũ sau bước vá thử. \textbf{Bằng chứng:} \texttt{data/rq1\_audit.csv},
\texttt{data/rq1\_cf\_identity.json}, \texttt{src/flx/bip\_run.py}.

\subsection{SemFault-250 đã sửa}

\textbf{Khẳng định.} Bộ dòng lỗi chuẩn được kiểm chứng lại bằng thực thi (\GtConfirmedN{} xác nhận, \GtCorrectedN{} sửa vị
trí); dùng nó làm Acc@1 của SemLoc tăng từ \RqTwoOriginalAllSemlocAccOneMean\% lên \RqTwoRepairedAllSemlocAccOneMean\%.
\textbf{Giới hạn:} bộ phán định chấp nhận cả bản sửa bù trừ; các vị trí trên dòng trống đã bị loại, còn bản sửa bù trừ trên dòng
có mã chưa được phát hiện tự động.

\subsection{Mơ hồ dòng kề cạnh của Acc@1 cấp dòng}

\textbf{Khẳng định.} Với lỗi trải trên dòng điều kiện và câu lệnh được bảo vệ, Acc@1 đổi hơn hai mươi điểm theo quy ước
gán nhãn (SemLoc \TolSemlocAccOneExact\% $\to$ \TolSemlocAccOneGuardStrict\%); thứ tự giữa các phương pháp giữ nguyên.
\textbf{Bằng chứng:} \texttt{data/failure\_tolerance.csv}.

\subsection{CovGate, bảng kiểm tra đánh giá và gói tái hiện}

CovGate tăng Acc@3 \RqTwoRepairedAllCovgateDAccThree{} điểm [\RqTwoRepairedAllCovgateDAccThreeLo;
\RqTwoRepairedAllCovgateDAccThreeHi] và giữ nguyên Acc@1 --- gắn thêm bằng chứng vào quy trình nhiều bước chỉ giúp ít so với đưa thẳng
bằng chứng cho LLM; với $\lambda$ chọn trên tập dev, CovGate thực chất là một bộ phá hòa bằng Ochiai. Bảng kiểm tra bảy mục
(phá hòa, mẫu số, dòng lỗi kiểm chứng, loại bỏ thành phần ẩn, lặp và kiểm định, chi phí, so với một prompt có cùng bằng
chứng) được áp cho mọi phương pháp. Gói tái hiện gồm Docker chạy không cần mạng, toàn bộ nhật ký và bộ đệm lượt
gọi LLM để phát lại, và macro sinh tự động cho mọi con số trong báo cáo.

\section{Triết lý khoa học}
\label{sec:triet-ly}

\subsection{Bằng chứng thực thi là trọng tài}

Ở mọi chỗ cần phán xử --- dòng nào là lỗi, bản sửa nào đúng --- nhóm để kiểm thử quyết định, LLM chỉ sinh ứng viên. Đây
cũng chính là câu hỏi đặt cho SemLoc: nếu một phương pháp tuyên bố neo lập luận vào thực thi, phần neo đó phải được đo.
Tái hiện không phải là chép lại con số: khớp số công bố là điều kiện cần; điều kiện đủ là hiểu vì sao ra con số đó và kết
quả còn đứng vững khi đổi cách chấm hợp lý.

\subsection{Đo đóng góp, so với phương pháp đối chứng đơn giản, và chấp nhận bị bác bỏ}

Một con số Acc@1 không nói nó đến từ đâu; phân rã đóng góp và loại bỏ từng thành phần biến câu hỏi ``phương pháp có tốt
không'' thành ``phần nào làm nó tốt''. Một quy trình nhiều bước chỉ đáng giá khi thắng phương pháp đối chứng đơn giản nhất có cùng
thông tin. Mỗi khẳng định đi kèm điều kiện bác bỏ và khoảng tin cậy; đơn vị lấy mẫu là chương trình (hoặc cụm khuôn
mẫu), không coi các lần chạy là mẫu độc lập; kết quả được báo cáo cả khi khoảng tin cậy chứa 0. Cuối cùng, bộ dữ liệu
cũng là một giả thuyết: nhãn do người gán có thể sai hoặc mơ hồ, và chính bộ phán định sửa nhãn cũng có giới hạn cần nêu rõ.

\section{Các mối đe dọa đến tính hợp lệ}
\label{sec:moi-de-doa}

\subsection{Tính hợp lệ nội tại và tính hợp lệ của cấu trúc}

\begin{itemize}
\item \textbf{Bộ phán định sửa tối thiểu} chấp nhận bản sửa bù trừ nên bộ dòng lỗi đã sửa còn nhiễu; mọi kết quả cũng được báo
cáo trên bộ gốc.
\item \textbf{Bộ dòng lỗi kiểm chứng có thể thiên về phương pháp dùng độ phủ}, vì bộ phán định cần kiểm thử thất bại chạy tới
dòng sửa; so sánh giữa SemLoc, các phương pháp đối chứng dùng LLM và CovGate cùng chiều trên bộ gốc, riêng cổng độ phủ thì không.
\item \textbf{Mơ hồ dòng kề cạnh}: Acc@1 chặt phụ thuộc quy ước gán nhãn; nhóm báo cáo cả Acc@1 chấp nhận đầu khối/thân khối.
\item \textbf{Prompt phương pháp đối chứng} do nhóm viết trước khi xem kết quả và không tinh chỉnh; prompt khác có thể cho kết quả khác.
\end{itemize}

\subsection{Tính hợp lệ ngoại tại}

\begin{itemize}
\item \textbf{Lỗi thật ít và chỉ một dự án} (31 lỗi youtube-dl); hàm đích được chọn từ bản sửa của lập trình viên nên định
vị cấp hàm không được đánh giá; một lỗi còn hai kiểm thử thất bại do môi trường trong bộ chạy gốc của SemLoc.
\item \textbf{Mô hình}: mô hình chính mới hơn mô hình trong bài gốc; mỗi nhà cung cấp chạy ở mức suy luận mặc định, nên
so sánh giữa họ mô hình là ``ở mặc định'' chứ không cùng mức suy luận; Claude qua cổng trung gian sinh ràng buộc yếu hơn
nhật ký công bố.
\item \textbf{Nhiễm dữ liệu}: biến thể đổi tên chỉ đổi định danh và xóa chú thích, cấu trúc mã giữ nguyên, nên không loại
trừ hoàn toàn việc mô hình đã thấy các khuôn mẫu.
\end{itemize}
